\subsubsection{具体分析}

问题一要求在稳定经济参数下给出逐地块七年种植计划，属于带跨期轮作约束的优化问题。本问采用确定性多期混合整数线性规划进行求解。

首先整理54个地块、41种作物和2023年种植统计，接着构造销量代理和价格代表值，然后建立面积、产销及轮作约束，最后分别求解滞销与半价销售模式并填写结果表。具体流程如图\ref{fig:analysis1_flow}所示。

\begin{figure}[H]
  \begin{center}
    \begin{tikzpicture}[x=1cm, y=0.7cm]
      \node[box] (n11) at (0, 0) {读取附件};
      \node[box] (n12) at (5, 0) {清洗数据};
      \node[box] (n13) at (10, 0) {构造销量};
      \node[box] (n23) at (10, -3) {建立变量};
      \node[box] (n22) at (5, -3) {添加约束};
      \node[box] (n21) at (0, -3) {两模式求解};
      \node[box] (n31) at (0, -6) {复算收益};
      \node[box] (n32) at (5, -6) {核验约束};
      \node[box] (n33) at (10, -6) {填写表格};
      \draw[arrow] (n11) -- (n12);
      \draw[arrow] (n12) -- (n13);
      \draw[arrow] (n13) -- ++(0,-1.0) -| (n23);
      \draw[arrow] (n23) -- (n22);
      \draw[arrow] (n22) -- (n21);
      \draw[arrow] (n21) -- (n31);
      \draw[arrow] (n31) -- (n32);
      \draw[arrow] (n32) -- (n33);
    \end{tikzpicture}
    \caption{问题一分析流程图}
    \label{fig:analysis1_flow}
  \end{center}
\end{figure}

\subsubsection{模型准备}

原始数据来自附件1和附件2。附件1含54个地块及41种作物的适宜性，附件2含87条2023年种植记录与107条亩产量、成本和价格记录。检查发现表内存在合并单元格，价格以区间字符串给出，智慧大棚第一季的部分参数需按附件说明派生，且附件没有独立的预期销量字段。

预处理分三步完成。第一步填充合并单元格锚点并统一地块、季次和作物编号。第二步解析价格区间并取中点作为确定性代表值。第三步用2023年种植面积乘对应亩产量，汇总得到销量代理$D_j^0$，并补齐智慧大棚第一季18个合法参数组合。处理后得到125个可用的地块类型、季次和作物参数组合。

该问题的目标和约束均可线性化。精确MILP能够输出完整面积方案和最优间隙，规则轮作可完成同一任务并作为基线，滚动时域仅在完整模型无法获得合格可行解时启用。基于人工确认的方法选择，本文使用精确MILP作为主方法\upcite{tan2025review,xu2025lp}。

\begin{longtable}{>{\centering\arraybackslash}p{0.20\textwidth}>{\centering\arraybackslash}p{0.22\textwidth}>{\centering\arraybackslash}p{0.22\textwidth}>{\centering\arraybackslash}p{0.22\textwidth}}
  \caption{问题一方法能力比较}\label{tab:q1_method_compare}\\
  \toprule
  方法 & 完整面积输出 & 跨年约束 & 求解质量证据 \\
  \midrule
  \endfirsthead
  \caption{问题一方法能力比较（续）}\\
  \toprule
  方法 & 完整面积输出 & 跨年约束 & 求解质量证据 \\
  \midrule
  \endhead
  \bottomrule
  \endfoot
  精确MILP & 支持 & 统一建模 & MIP Gap \\
  规则轮作基线 & 支持 & 显式构造 & 可行性核验 \\
  滚动时域备用 & 支持 & 需保留边界状态 & 分阶段Gap \\
\end{longtable}
